\uuid{geFO}
\titre{Limite en un point}
\auteur{bodin ; hanani ; mzari}
\organisation{exo7}
\licence{CC BY-NC-SA 4.0}
\datecreate{2020-01-20}
\niveau{L1}
\module{Analyse}
\chapitre{Continuité, limite et étude de fonctions réelles}
\sousChapitre{Limite de fonctions}
\difficulte{3}
\theme{Limites des fonctions réelles, limite en un point}
\exo7qcmid{lille-1.limites.021}

\contenu{
\begin{question}
\qclassification{123.03}
\qtags{motcle=limite en un point}
Soit $f(x)= xE(\frac{1}{x})$, où $E$ désigne la partie entière. Quelles sont les assertions vraies ?
\begin{answers}

    \bad{$\lim_{x\to 0} f(x)=0$}

    \bad{$\lim_{x\to 0} f(x)=+\infty$}

    \bad{$f$ n'admet pas de limite en $0$.}

    \good{$\lim_{x\to 0} f(x)=1$}
\end{answers}
\begin{explanations}
Pour tout $x\in \Rr$, on a : $x-1<E(x)\le x$. Donc $1-x < f(x) \le 1$, pour $x>0$ et
$1 \le f(x) < 1-x$, pour $x<0$. On déduit que $\lim_{x\to 0} f(x) =1$.
\end{explanations}
\end{question}
}
